\emph{Relasi}~$R$ ing himpunan~$A$ nggandhèngaké anggota~$A$.
Kita nulis $Rxy$ yèn $x$ lan~$y$ \emph{gegayutan} miturut~$R$.
Sacara formal, $R$ iku himpunan pasangan $\tuple{x,y}\in A^2$
kang nyukupi~$Rxy$. Tuladha wigati yaiku luwih cilik, luwih gedhé,
padha, mbagi tanpa turahan, padha dawané, dadi subhimpunan, lan
padha gedhéné. Tuladha kasebut bisa gegayutan karo wilangan,
rerangken, utawa himpunan. \emph{Graf} minangka cara umum kanggo
nggambar relasi. Graf dhéwé uga bisa dianggep minangka relasi
biner, yaiku relasi \emph{rusuk}, bebarengan karo himpunan \emph{simpul}.

Sawetara relasi nduwèni sipat wigati. Relasi~$R$ iku \emph{refleksif}
yèn saben unsur nduwèni relasi~$R$ karo awaké dhéwé, \emph{simètris}
yèn $Rxy$ njalari $Ryx$ kanggo sembarang $x$ lan~$y$, lan
\emph{transitif} yèn $Rxy$ lan $Ryz$ njalari~$Rxz$. Relasi kang
nduwèni telung sipat kasebut diarani \emph{relasi ekuivalènsi}.
Relasi iku \emph{antisimètris} yèn $Rxy$ lan $Ryx$ njalari $x=y$.
\emph{Urutan parsial} iku relasi kang refleksif, antisimètris, lan
transitif. \emph{Urutan linier} iku urutan parsial kang, kanggo
sembarang $x$ lan~$y$, nyukupi $Rxy$, $x=y$, utawa $Ryx$.
Sacara umum relasi kang nduwèni sipat pungkasan iki diarani
\emph{gegayutan}.

Amarga relasi iku himpunan pasangan, operasi himpunan kayata
gabungan lan irisan bisa ditrapaké marang relasi. Relasi uga
bisa diranté nganggo \emph{prodhuk relatif} $R\mid S$.
Gabungan prodhuk relatif saka~$R$ karo awaké dhéwé kanggo
kabèh cacah langkah positif ngasilaké \emph{panutupan transitif}
$R^+$ saka~$R$.
